\newpage % Rozdziały zaczynamy od nowej strony.
\section{Wyniki ewaluacji eksperymentalnej}

Niniejszy rozdział przedstawia wyniki eksperymentalnej ewaluacji zaproponowanego frameworku do predykcji mutacji aminokwasowych w białkach wirusowych. Celem przeprowadzonych eksperymentów było zbadanie skuteczności trzech architektur modeli opartych na sieciach LSTM w~zadaniu klasyfikacji binarnej pozycji aminokwasowych jako ulegających bądź nieulegających mutacji. Ewaluację przeprowadzono na dwóch zbiorach danych --- sekwencjach białka S wirusa SARS-CoV-2 oraz białka hemagglutyniny (HA) wirusa Influenza A podtypu H1N1 --- co pozwoliło zweryfikować zdolność frameworku do generalizacji między różnymi patogenami wirusowymi.

%--------------------------------------
\subsection{Środowisko eksperymentalne}
%--------------------------------------

Wszystkie eksperymenty przeprowadzono w~zunifikowanym środowisku obliczeniowym w~celu zapewnienia powtarzalności wyników. Konfiguracja sprzętowa i~programowa została zestawiona w~tabeli~\ref{tab:env}.

\begin{table}[H]
    \centering
    \caption{Konfiguracja środowiska eksperymentalnego.}
    \label{tab:env}
    \begin{tabular}{ll}
        \toprule
        \textbf{Parametr} & \textbf{Wartość} \\
        \midrule
        Procesor graficzny     & % TODO: np. NVIDIA RTX 3090 / Apple M1 Pro \\
        Pamięć RAM             & % TODO: np. 32 GB \\
        System operacyjny      & Ubuntu 20.04 (Docker) \\
        Język programowania    & Python 3.x \\
        Framework ML           & PyTorch~\cite{paszke2019pytorch} \\
        Biblioteka ML          & scikit-learn~\cite{pedregosa2011sklearn} \\
        Ziarno losowości       & 42 \\
        Format konfiguracji    & YAML \\
        \bottomrule
    \end{tabular}
\end{table}

W~celu zapewnienia reprodukowalności wyników zastosowano stałe ziarno losowości (\textit{random seed}) o~wartości 42, które inicjalizowało generatory liczb pseudolosowych w~bibliotekach PyTorch, NumPy oraz wbudowanym module \texttt{random} języka Python. Konfiguracja każdego eksperymentu została zapisana w~pliku YAML, co umożliwia dokładne odtworzenie warunków eksperymentalnych.

%--------------------------------------
\subsection{Zbiory danych}
%--------------------------------------

Ewaluację przeprowadzono na dwóch zbiorach danych, różniących się źródłem, rozmiarem oraz granularnością podziału czasowego. Szczegółowe porównanie parametrów obu zbiorów przedstawiono w~tabeli~\ref{tab:datasets}.

\begin{table}[H]
    \centering
    \caption{Porównanie zbiorów danych wykorzystanych w~eksperymentach.}
    \label{tab:datasets}
    \begin{tabular}{lcc}
        \toprule
        \textbf{Parametr} & \textbf{SARS-CoV-2 Spike} & \textbf{Influenza A H1N1 HA} \\
        \midrule
        Źródło danych          & GISAID~\cite{elbe2017gisaid}  & NCBI~\cite{bao2008ncbi} \\
        Białko                 & Spike (S)                     & Hemagglutynina (HA) \\
        Długość sekwencji      & 1273 aa                       & 565 aa \\
        Liczba epitopów        & 16 regionów (317 poz.)        & 5 miejsc (95 poz.) \\
        Granularność czasowa   & Miesięczna                    & Roczna \\
        Rozmiar zbioru         & $\sim$12\,680 próbek          & $\sim$40\,000 próbek \\
        Okno czasowe           & 10 miesięcy                   & 10 lat \\
        Stosunek mutacji       & $\sim$20\%                    & $\sim$20\% \\
        Podział train/val/test & 80\%\,/\,10\%\,/\,10\%       & 80\%\,/\,10\%\,/\,10\% \\
        \bottomrule
    \end{tabular}
\end{table}

Oba zbiory danych zostały przygotowane z~wykorzystaniem tego samego potoku przetwarzania, obejmującego: filtrowanie sekwencji, kodowanie metodą ProtVec~\cite{asgari2015protvec}, klasteryzację algorytmem K-means~\cite{macqueen1967kmeans} w~ramach poszczególnych okresów czasowych, łączenie klastrów między kolejnymi okresami na podstawie odległości centroidów oraz generowanie próbek treningowych przy użyciu okna przesuwnego. Podział na zbiór treningowy, walidacyjny i~testowy został przeprowadzony chronologicznie --- zbiór testowy obejmuje najnowsze okresy czasowe, co odzwierciedla rzeczywisty scenariusz predykcji przyszłych mutacji.

%--------------------------------------
\subsection{Metryki ewaluacji}
%--------------------------------------

Ze względu na niezbalansowany charakter zadania klasyfikacji binarnej (stosunek klas mutacja/brak mutacji wynosi około 20\%/80\%) zastosowano zestaw metryk uwzględniających tę asymetrię. Oprócz standardowej dokładności (ang.~\textit{accuracy}) wykorzystano następujące miary:

\begin{itemize}
    \item \textbf{Precyzja} (ang.~\textit{precision}) --- stosunek prawidłowo przewidzianych mutacji do~wszystkich przewidzianych mutacji:
    \begin{equation}
        \text{Precision} = \frac{TP}{TP + FP}
        \label{eq:precision}
    \end{equation}

    \item \textbf{Czułość} (ang.~\textit{recall}) --- stosunek prawidłowo przewidzianych mutacji do~wszystkich rzeczywistych mutacji:
    \begin{equation}
        \text{Recall} = \frac{TP}{TP + FN}
        \label{eq:recall}
    \end{equation}

    \item \textbf{Miara F1} (ang.~\textit{F1 score}) --- średnia harmoniczna precyzji i~czułości:
    \begin{equation}
        F_1 = 2 \cdot \frac{\text{Precision} \cdot \text{Recall}}{\text{Precision} + \text{Recall}}
        \label{eq:f1}
    \end{equation}

    \item \textbf{Współczynnik korelacji Matthewsa} (ang.~\textit{Matthews Correlation Coefficient}, MCC)~\cite{matthews1975mcc} --- miara jakości klasyfikacji binarnej uwzględniająca wszystkie elementy macierzy pomyłek:
    \begin{equation}
        \text{MCC} = \frac{TP \cdot TN - FP \cdot FN}{\sqrt{(TP+FP)(TP+FN)(TN+FP)(TN+FN)}}
        \label{eq:mcc-eval}
    \end{equation}

    \item \textbf{Pole pod krzywą ROC} (ang.~\textit{ROC-AUC}) --- miara zdolności modelu do~rozróżniania klas niezależnie od~progu decyzyjnego.
\end{itemize}

Współczynnik MCC przyjmuje wartości z~przedziału $[-1, 1]$, gdzie $1$ oznacza idealną klasyfikację, $0$ --- klasyfikację na poziomie losowym, a~$-1$ --- klasyfikację całkowicie odwrotną. Jak wykazali Chicco i~Jurman~\cite{chicco2020mcc}, MCC stanowi bardziej wiarygodną miarę jakości klasyfikacji binarnej niż dokładność czy miara F1 w~przypadku zbiorów niezbalansowanych, dlatego metryka ta została przyjęta jako główne kryterium porównawcze modeli.

%--------------------------------------
\subsection{Model bazowy}
%--------------------------------------

Jako punkt odniesienia (ang.~\textit{baseline}) dla modeli opartych na sieciach LSTM zastosowano klasyfikator regresji logistycznej (ang.~\textit{Logistic Regression}), zaimplementowany z~wykorzystaniem biblioteki scikit-learn~\cite{pedregosa2011sklearn}. Model bazowy otrzymywał na wejściu spłaszczoną reprezentację okna czasowego --- konkatenację wektorów ProtVec ze~wszystkich kroków czasowych w~oknie --- i~dokonywał klasyfikacji binarnej każdej pozycji aminokwasowej niezależnie. Wybór regresji logistycznej jako modelu bazowego pozwala ocenić, w~jakim stopniu architektury rekurencyjne wykorzystują informację o~zależnościach temporalnych w~porównaniu z~modelem liniowym nieuwzględniającym struktury sekwencyjnej danych.

%--------------------------------------
\subsection{Wyniki dla SARS-CoV-2 Spike}
\label{sec:results-sarscov2}
%--------------------------------------

\subsubsection{Porównanie architektur}

W~tabeli~\ref{tab:results-sarscov2} zestawiono wyniki ewaluacji czterech modeli na zbiorze testowym sekwencji białka S wirusa SARS-CoV-2. Wartości metryk zostały obliczone na zbiorze testowym obejmującym chronologicznie najnowsze okresy czasowe.

\begin{table}[H]
    \centering
    \caption{Porównanie wyników modeli na zbiorze testowym SARS-CoV-2 Spike. Najlepsze wartości dla każdej metryki zostały wyróżnione pogrubieniem.}
    \label{tab:results-sarscov2}
    \begin{tabular}{lcccccc}
        \toprule
        \textbf{Model} & \textbf{Accuracy} & \textbf{Precision} & \textbf{Recall} & \textbf{F1} & \textbf{MCC} & \textbf{ROC-AUC} \\
        \midrule
        Regresja logistyczna    & --- & --- & --- & --- & --- & --- \\
        LSTM                    & --- & --- & --- & --- & --- & --- \\
        Attention LSTM          & --- & --- & --- & --- & --- & --- \\
        Dual Attention LSTM     & --- & --- & --- & --- & --- & --- \\
        \bottomrule
    \end{tabular}
    % TODO: Uzupełnić wartości metryk po przeprowadzeniu eksperymentów.
\end{table}

% TODO: Analiza wyników --- który model osiągnął najlepsze wyniki i z jakiego powodu.
% Oczekiwania: modele z mechanizmem atencji powinny osiągać wyższe wartości MCC
% niż standardowy LSTM, ponieważ mogą selektywnie ważyć informację z różnych
% kroków czasowych okna. Dual Attention LSTM, operujący zarówno na wymiarze
% czasowym, jak i wymiarze cech, powinien osiągnąć najwyższe wartości.

% TODO: Rysunek - krzywe ROC dla wszystkich modeli na jednym wykresie
% \begin{figure}[H]
%     \centering
%     % \includegraphics[width=0.7\textwidth]{rysunki/roc_sarscov2.pdf}
%     \caption{Krzywe ROC dla porównywanych modeli na zbiorze testowym SARS-CoV-2 Spike.}
%     \label{fig:roc-sarscov2}
% \end{figure}

% TODO: Rysunek - macierz pomyłek najlepszego modelu
% \begin{figure}[H]
%     \centering
%     % \includegraphics[width=0.5\textwidth]{rysunki/confusion_matrix_sarscov2.pdf}
%     \caption{Macierz pomyłek najlepszego modelu na zbiorze testowym SARS-CoV-2 Spike.}
%     \label{fig:cm-sarscov2}
% \end{figure}

\subsubsection{Analiza mechanizmu atencji}

Istotną zaletą modeli wyposażonych w~mechanizm atencji jest możliwość interpretacji nauczonych wag, które wskazują, którym krokom czasowym w~oknie przesuwnym model przypisuje największe znaczenie predykcyjne.

% TODO: Rysunek - wizualizacja wag atencji (heatmapa lub wykres słupkowy)
% pokazujący średnie wagi atencji dla poszczególnych kroków czasowych w oknie.
% \begin{figure}[H]
%     \centering
%     % \includegraphics[width=0.7\textwidth]{rysunki/attention_weights_sarscov2.pdf}
%     \caption{Średnie wagi atencji modelu Attention LSTM dla poszczególnych kroków
%              czasowych w oknie przesuwnym (SARS-CoV-2 Spike).}
%     \label{fig:attention-sarscov2}
% \end{figure}

% TODO: Interpretacja biologiczna wag atencji:
% - Czy model przypisuje większą wagę nowszym okresom (bliższym predykowanemu)?
% - Czy obserwuje się wzrost wag w okresach pojawienia się wariantów Alpha, Delta, Omicron?
% - Jakie wnioski biologiczne można wyciągnąć z rozkładu wag atencji?

%--------------------------------------
\subsection{Wyniki dla Influenza A H1N1 HA}
\label{sec:results-influenza}
%--------------------------------------

\subsubsection{Porównanie architektur}

Analogiczną ewaluację przeprowadzono na zbiorze danych sekwencji białka HA wirusa Influenza A podtypu H1N1. Wyniki zestawiono w~tabeli~\ref{tab:results-influenza}.

\begin{table}[H]
    \centering
    \caption{Porównanie wyników modeli na zbiorze testowym Influenza A H1N1 HA. Najlepsze wartości dla każdej metryki zostały wyróżnione pogrubieniem.}
    \label{tab:results-influenza}
    \begin{tabular}{lcccccc}
        \toprule
        \textbf{Model} & \textbf{Accuracy} & \textbf{Precision} & \textbf{Recall} & \textbf{F1} & \textbf{MCC} & \textbf{ROC-AUC} \\
        \midrule
        Regresja logistyczna    & --- & --- & --- & --- & --- & --- \\
        LSTM                    & --- & --- & --- & --- & --- & --- \\
        Attention LSTM          & --- & --- & --- & --- & --- & --- \\
        Dual Attention LSTM     & --- & --- & --- & --- & --- & --- \\
        \bottomrule
    \end{tabular}
    % TODO: Uzupełnić wartości metryk po przeprowadzeniu eksperymentów.
\end{table}

% TODO: Analiza wyników dla Influenza --- porównanie modeli, interpretacja.
% Dane Influenza mają roczną granularność i dłuższą historię ewolucyjną,
% co może wpływać na jakość predykcji w porównaniu z danymi SARS-CoV-2.

\subsubsection{Analiza mechanizmu atencji}

% TODO: Wizualizacja i interpretacja wag atencji dla danych Influenza.
% - Czy rozkład wag różni się od SARS-CoV-2?
% - Jak roczna granularność wpływa na rozkład wag?
% - Czy model identyfikuje lata o szczególnym znaczeniu epidemiologicznym?

% \begin{figure}[H]
%     \centering
%     % \includegraphics[width=0.7\textwidth]{rysunki/attention_weights_influenza.pdf}
%     \caption{Średnie wagi atencji modelu Attention LSTM dla poszczególnych kroków
%              czasowych w oknie przesuwnym (Influenza A H1N1 HA).}
%     \label{fig:attention-influenza}
% \end{figure}

%--------------------------------------
\subsection{Analiza porównawcza}
\label{sec:cross-analysis}
%--------------------------------------

W~celu oceny zdolności frameworku do generalizacji między różnymi białkami wirusowymi przeprowadzono analizę porównawczą wyników uzyskanych dla obu zbiorów danych.

% TODO: Tabela podsumowująca - najlepsze wyniki MCC dla każdego modelu i zbioru danych.
% \begin{table}[H]
%     \centering
%     \caption{Porównanie najlepszych wartości MCC między zbiorami danych.}
%     \label{tab:cross-comparison}
%     \begin{tabular}{lcc}
%         \toprule
%         \textbf{Model} & \textbf{SARS-CoV-2 Spike} & \textbf{Influenza A H1N1 HA} \\
%         \midrule
%         Regresja logistyczna    & --- & --- \\
%         LSTM                    & --- & --- \\
%         Attention LSTM          & --- & --- \\
%         Dual Attention LSTM     & --- & --- \\
%         \bottomrule
%     \end{tabular}
% \end{table}

% TODO: Analiza porównawcza obejmująca następujące aspekty:
% 1. Czy ranking modeli jest spójny między oboma zbiorami danych?
%    Jeśli tak, potwierdza to, że przewaga architektur z atencją jest niezależna
%    od konkretnego białka wirusowego.
% 2. Wpływ rozmiaru zbioru danych na jakość predykcji --- zbiór Influenza
%    jest ponad trzykrotnie większy niż zbiór SARS-CoV-2.
% 3. Wpływ granularności czasowej --- miesięczna (SARS-CoV-2) vs. roczna (Influenza).
% 4. Różnice w dynamice ewolucyjnej obu wirusów i ich przełożenie na wyniki.

%--------------------------------------
\subsection{Dyskusja}
\label{sec:discussion}
%--------------------------------------

% TODO: Interpretacja wyników w kontekście literatury.
% Porównanie z podejściami opartymi na modelach językowych (Hie i in.~\cite{hie2021viral}),
% które wykorzystywały wstępnie wytrenowane reprezentacje sekwencji białkowych
% do przewidywania ucieczki immunologicznej wirusów. W odróżnieniu od tego podejścia,
% zaproponowany framework jawnie modeluje dynamikę temporalną ewolucji wirusowej
% z wykorzystaniem okna przesuwnego i architektur rekurencyjnych.

% TODO: Porównanie z podejściem Mahera i in.~\cite{maher2022predicting},
% którzy zastosowali modele gradient boosting do identyfikacji sterowników
% mutacyjnych wariantów SARS-CoV-2. Zaproponowany framework różni się
% wykorzystaniem informacji temporalnej z wielu kolejnych okresów czasowych
% oraz możliwością zastosowania do różnych białek wirusowych.

% TODO: Omówienie ograniczeń eksperymentów:
% - Zależność od jakości danych źródłowych (bias w próbkowaniu sekwencji).
% - Wpływ progów klasteryzacji na jakość predykcji.
% - Ograniczenia wynikające z binarnej natury predykcji (mutacja/brak mutacji)
%   w porównaniu z predykcją konkretnego aminokwasu substytucyjnego.
% - Brak walidacji prospektywnej na danych z przyszłych okresów.

% TODO: Wnioski z eksperymentów:
% - Podsumowanie, która architektura okazała się najskuteczniejsza.
% - Ocena, czy mechanizm atencji wnosi istotną poprawę.
% - Ocena konfigurowalności i reprodukowalności frameworku.
% - Wnioski dotyczące generalizowalności między różnymi białkami wirusowymi.
